La auditoría por \texttt{discoverymethod} identifica nodos y componentes dominados por ciertos métodos de descubrimiento. Sin embargo, esa observación por sí sola no basta: si un método es muy frecuente en el catálogo global, algunos nodos podrían parecer dominados por ese método simplemente por azar. Por esta razón, se implementó una prueba nula por permutación.

La hipótesis nula de esta prueba es que las etiquetas \texttt{discoverymethod} no están asociadas de forma especial con la estructura Mapper. En otras palabras, si los métodos de descubrimiento se repartieran aleatoriamente entre los mismos planetas, los nodos no deberían mostrar una concentración por método mayor que la esperada bajo azar.

La prueba mantiene fija toda la estructura de Mapper: la topología del grafo, los nodos, las aristas y la membresía de planetas en cada nodo. El Mapper no se vuelve a correr. Lo único que se permuta son las etiquetas \texttt{discoverymethod} entre planetas. Esto permite construir escenarios nulos donde los mismos nodos existen, con los mismos planetas y la misma distribución global de métodos, pero sin asociación específica entre método de descubrimiento y posición topológica.

Esta decisión es importante. Si se permutaran las variables físicas u orbitales, cambiaría la geometría del problema y se estaría evaluando otra cosa. Aquí la pregunta es más precisa: dado el grafo Mapper observado, ¿la concentración de métodos de descubrimiento dentro de sus nodos es mayor que la esperada por azar?

Para cada nodo \(v\), se compara la pureza observada por método con la distribución nula obtenida por permutación. El estadístico reportado es:

\[
Z_v =
\frac{
P_v^{obs}-\mathbb{E}[P_v^{null}]
}{
\operatorname{sd}(P_v^{null})+\epsilon
}.
\]

Aquí \(P_v^{obs}\) es la pureza observada del nodo, \(\mathbb{E}[P_v^{null}]\) es la pureza promedio esperada bajo permutaciones, \(\operatorname{sd}(P_v^{null})\) es su desviación estándar y \(\epsilon\) es un término pequeño para evitar división entre cero. Un valor \(Z_v\) alto indica que la concentración observada por método es mucho mayor que la esperada bajo la hipótesis nula.

Además de los enriquecimientos por nodo, se calcularon estadísticos globales por grafo. Los principales fueron: la pureza media ponderada por tamaño de nodo, la entropía media ponderada y la información mutua normalizada, o NMI, entre una asignación dura de planeta a nodo y \texttt{discoverymethod}.

NMI se interpreta aquí como una medida de cuánta información comparte la posición de un planeta en el grafo con su método de descubrimiento. Si el NMI observado es alto en comparación con la distribución nula, entonces saber en qué nodo cae un planeta ayuda a predecir con qué método fue descubierto. Como Mapper permite que un planeta pertenezca a varios nodos, se construyó una asignación dura determinista: cada planeta se asignó al nodo más grande al que pertenece; en caso de empate, se usó el orden alfabético del identificador del nodo.

La auditoría final se realizó con 1000 permutaciones. Por lo tanto, valores empíricos como \(p=0.001\) indican evidencia fuerte de asociación bajo la resolución de esta prueba. Esta evidencia no prueba causalidad, pero sí muestra que la alineación entre topología y método de descubrimiento es demasiado fuerte para atribuirla fácilmente al azar.

\begin{table}[H]
\centering
\scriptsize
\caption{Prueba nula por permutacion de etiquetas discoverymethod.}
\label{tab:bias_permutation_null}
\begin{adjustbox}{max width=\linewidth}
\begin{tabular}{lrrrrrrrr}
\toprule
config\_id & observed\_weighted\_mean\_method\_purity & null\_mean\_purity & purity\_z & purity\_empirical\_p & observed\_nmi & null\_mean\_nmi & nmi\_z & nmi\_empirical\_p \\
\midrule
phys\_min\_pca2\_cubes10\_overlap0p35 & 0.765 & 0.730 & 12.961 & 0.001 & 0.136 & 0.009 & 131.786 & 0.001 \\
phys\_density\_pca2\_cubes10\_overlap0p35 & 0.783 & 0.731 & 21.768 & 0.001 & 0.150 & 0.009 & 136.246 & 0.001 \\
orbital\_pca2\_cubes10\_overlap0p35 & 0.842 & 0.773 & 26.630 & 0.001 & 0.254 & 0.009 & 209.663 & 0.001 \\
joint\_no\_density\_pca2\_cubes10\_overlap0p35 & 0.851 & 0.748 & 37.830 & 0.001 & 0.226 & 0.009 & 233.752 & 0.001 \\
joint\_pca2\_cubes10\_overlap0p35 & 0.848 & 0.746 & 38.754 & 0.001 & 0.221 & 0.009 & 232.721 & 0.001 \\
thermal\_pca2\_cubes10\_overlap0p35 & 0.813 & 0.756 & 20.860 & 0.001 & 0.207 & 0.012 & 147.018 & 0.001 \\
\bottomrule
\end{tabular}
\end{adjustbox}

\end{table}


La Tabla anterior resume los resultados globales de la prueba por permutación para los grafos seleccionados. El caso más importante para la interpretación del reporte es \texttt{orbital\_pca2\_cubes10\_overlap0p35}, porque este grafo fue identificado como el candidato estructural principal por su alta complejidad topológica y baja imputación.

Para el grafo orbital, la pureza media ponderada observada por método fue 0.842, mientras que la media nula de pureza fue 0.773. El valor \texttt{purity\_z} fue 26.630 y el valor empírico fue \(p=0.001\). Esto indica que los nodos del grafo orbital están más concentrados por método de descubrimiento de lo esperado bajo permutación aleatoria.

El resultado es todavía más claro en NMI. El NMI observado para el grafo orbital fue 0.254, mientras que el NMI nulo medio fue 0.009. El valor \texttt{nmi\_z} fue 209.663 y el valor empírico fue \(p=0.001\). En términos interpretativos, esto significa que la asignación de planetas a nodos en el grafo orbital contiene información significativa sobre \texttt{discoverymethod}. Es decir, la topología orbital no es independiente del proceso de descubrimiento.

Este resultado debe leerse con cuidado. No significa que el grafo orbital sea únicamente un artefacto observacional. El mismo grafo tiene \(\beta_1=96\), imputación media nodal de 0.011 y está construido con variables físicamente relevantes, \texttt{pl\_orbper} y \texttt{pl\_orbsmax}. Por lo tanto, sigue siendo el candidato más informativo para inspección científica. Sin embargo, la prueba de permutación muestra que su estructura no puede interpretarse como señal orbital pura.

La conclusión adecuada es mixta. El Mapper orbital revela una estructura rica y de baja imputación, compatible con organización por escala orbital, pero esa estructura también está fuertemente asociada con el método de descubrimiento. Esto es astrofísicamente razonable: los métodos de tránsito favorecen planetas cercanos y de periodo corto, mientras que otros métodos, como velocidad radial o imagen directa, favorecen distintas regiones del espacio orbital y físico. Por tanto, parte de la topología orbital puede reflejar arquitectura planetaria y parte puede reflejar selección observacional.

Esta prueba nula cumple una función central en el reporte. Sin ella, una región dominada por planetas orbitalmente parecidos podría interpretarse demasiado rápido como una población física. Con la permutación, en cambio, se puede distinguir entre regiones con señal más limpia y regiones donde la topología está significativamente alineada con \texttt{discoverymethod}. Por ello, los resultados de esta sección alimentan directamente la síntesis final de regiones: una región con coherencia física u orbital pero enriquecimiento significativo por método debe clasificarse como \texttt{mixed}, no como puramente física.

En resumen, la prueba por permutación muestra que la topología Mapper contiene asociación medible con el proceso de descubrimiento. Esta asociación no invalida el análisis, pero sí limita el tipo de conclusión que puede formularse. El resultado principal no es que Mapper descubra una topología orbital pura, sino que identifica una estructura exploratoria donde se mezclan señal orbital, sesgo observacional y decisiones de procesamiento. Esa es precisamente la razón por la que el reporte adopta una clasificación final de regiones físicas, observacionales, mixtas o débiles.




